\section{Related Work}

\textbf{Temporal dynamics of EPSS.} Ravalico et al.~\cite{ravalico2025epsstemporal} examined the temporal stability of EPSS scores across approximately 45{,}000 CVEs and reported meaningful score drift over 30-day periods. That work studies the per-CVE temporal evolution of the EPSS score itself. The present paper complements it by studying a different temporal effect: how EPSS's \textit{prioritization effectiveness}, evaluated as a ranking signal under active remediation, degrades as the high-EPSS subset of the backlog is exhausted. The two effects are independent in principle and additive in practice; the prioritization-effectiveness collapse documented here would occur even with perfectly stable per-CVE EPSS scores.

\textbf{Risk-based vulnerability management.} Jacobs et al.~\cite{jacobs2020improving} demonstrated that exploit-prediction models substantially outperform CVSS-based triage on enterprise data. Their evaluation is single-period in the sense relevant to this paper: it does not measure how the comparative advantage of an exploit-prediction-driven policy changes as the policy's recommended actions are executed across consecutive windows.

\textbf{Patch-management scheduling.} Optimisation-based formulations of patch scheduling treat the full multi-period problem with explicit capacity, deadline, and dependency constraints. The present paper sits at a different level of abstraction: it evaluates greedy single-window scoring heuristics applied sequentially, which is the dominant practical approach in enterprise operations and the regime under which EPSS is most commonly deployed.

\textbf{Cyber-hygiene benchmarking.} HygieneBench (Paper~3)~\cite{paper3hygienebench} established that hygiene signals --- patch posture and AD group drift in particular --- carry discriminative structure in synthetic anomaly-detection tasks. Paper~4~\cite{paper4hygieneprio} (HygienePrio) operationalised those signals in a single-window prioritization scorer. This paper extends Paper~4 along the temporal axis, isolating the multi-window setting as the locus where the case for hygiene-based scoring is operationally strongest.

\textbf{Vulnerability scoring fundamentals.} CVSS v3.1~\cite{first2019cvss} provides a standardised severity framework but is asset-agnostic; EPSS~\cite{jacobs2021epss} estimates per-CVE exploit likelihood without regard to host state. The KEV catalog~\cite{cisa2026kev} flags actively exploited CVEs; recency-of-KEV inclusion is used as one component of HygienePrio (\S5.1).

\textbf{Policy mandates.} CISA BOD~22-01~\cite{cisa2021bod2201} and BOD~23-01~\cite{cisa2022bod2301}, together with NIST SP~800-40 Rev.~4~\cite{nist2022sp80040r4}, define the operational scheduling regime modelled here. To our knowledge, no prior work has reported an empirical multi-window prioritization study under a pre-registered evaluation protocol with frozen seeds and bias-corrected percentile bootstrap confidence intervals (an approximation to the pre-registered BCa).
